\documentclass[a4paper,10pt]{article}
\usepackage{paper-en}
\usepackage{hyperref}

\def\thetitle{(Sym)metric plane spotting}
\def\theauthors{Nina Lebedeva and Anton Petrunin}

\hypersetup{colorlinks=true,
citecolor=black,
linkcolor=black,
anchorcolor=black,
filecolor=black,
menucolor=black,
urlcolor=black,
pdftitle={\thetitle},
pdfauthor={\theauthors}
}

%\overfullrule=100mm

\begin{document}

\title{\thetitle}
\author{\theauthors}
\date{}
\maketitle

\begin{abstract}
We characterize Euclidean and Lobachevsky planes as metric spaces with three metric conditions: existence of geodesics, side-side-side congruence, and uniqueness of reflection.
\end{abstract}

We assume that the reader is not afraid of metric spaces or the axiomatic approach to classical geometry.

\begin{thm}{Theorem}\label{thm:main}
A nontrivial 3-point-homogeneous geodesic space is isometric to a Euclidean and Lobachevsky plane
if for any two distinct points there is exactly one nonidentity isometry fixing both.
\end{thm}

Lobachevsky plane can have arbitrary negative curvature;
in particular rescaling of a Lobachevsky plane is another Lobachevsky plane.

Let us spell out the terms.
A space is \emph{nontrivial} if it has at least two points.

A \emph{geodesic} in a metric space $M$ is an isometric map from a real interval to $M$.
The image of a geodesic with endpoints $x$ and $y$ will be denoted by $[xy]$; we also call it a \emph{geodesic}.
If any two points in $M$ can be joined by a geodesic, then $M$ is called a \emph{geodesic space}.
Thus the geodesic condition is a weakened form of Euclid's first postulate: \textit{a straight line can be drawn from any point to any point}.

An \emph{isometry} of $M$ is a bijection from $M$ to itself that preserves distances.
Three-point homogeneity is essentially the side-side-side congruence criterion.
A metric space $M$ is \emph{3-point-homogeneous} if for any points $x_1$, $x_2$, $x_3$, $y_1$, $y_2$, $y_3\in M$ the equalities $|x_i-x_j|=|y_i-y_j|$ for all $i$ and $j$ imply that some isometry of $M$ sends $x_i$ to $y_i$ for each~$i$.
Here $|p-q|$ denotes the distance between points $p,q\in M$.
The points in a triple need not be distinct; in particular, 3-point homogeneity implies 2-point homogeneity and ordinary (1-point) homogeneity.

In a Euclidean or Lobachevsky plane, the only nonidentity isometry fixing two distinct points is, of course, reflection in the line thru them.
Thus the last condition of the theorem can be regarded as the \emph{unique reflection axiom}.
To characterize the Euclidean plane, we can add the following \emph{doubling axiom}:

\parit{The space is isometric to a copy of itself with all distances doubled.}

\medskip

It follows from the theorem since Lobachevsky plane is not isometric to a copy of itself with doubled metric.

\medskip

Note that we write our axioms in purely metric terms.
This naturally brings together metric geometry and the foundations of geometry.
A closely related but not purely metric approach was proposed by George Birkhoff \cite{birkhoff1932}.
Using side-side-side congruence in this setting is not new.
It was used by Karol Borsuk and Wanda Szmielew \cite[part I chap. II § 4]{borsuk-szmielew}, and by Alexandr Alexandrov \cite{alexandrov1994};
in both cases, equality of angles is defined thru equality of sides of auxiliary triangles.

Note that we do not assume local compactness; otherwise the theorem would follow at once from a theorem of Jacques Tits~\cite[IV.E § 3.59 cor. 2]{tits}, which classifies locally compact 2-point-homogeneous geodesic spaces.
All the ideas in the proof are elementary, and many are ancient.
This note was prompted by another result of ours~\cite{lebedeva-petrunin}.%%%???

\section*{Proof}

Assume that a metric space $M$ satisfies the assumptions of our theorem.

\begin{wrapfigure}{r}{35mm}
\centering
\vskip-5mm
\includegraphics{mppics/pic-10}
\vskip-0mm
\end{wrapfigure}

\parbf{Two distances.}
Let $p,q\in M$ be two distinct points.
The distances $r=|p-x|$ and $s=|q-x|$ leave at most two possibilities for $x$.

Indeed, there is a unique nonidentity isometry, say $\sigma$, fixing $p$ and~$q$.
If $|p-y|=r$ and $|q-y|=s$, then 3-point homogeneity gives an isometry fixing $p$ and $q$ and sending $x$ to $y$.
Thus $y=x$ if the isometry is the identity, and $y=\sigma(x)$ otherwise --- there is no third option.

\parbf{Local uniqueness of geodesics.}
We will show that there is $\eps_0>0$ such that any two points $x$ and $y$ at distance less than $\eps_0$ are joined by a unique geodesic.

A \emph{midpoint} of $x$ and $y$ is a point $m$ such that
\[|x-m|=|y-m|=\tfrac12\cdot|x-y|.\]
If the midpoint is unique, then every geodesic $[xy]$ contains it.
If midpoints are unique for all pairs of points at distance less than $\eps_0$, then repeated bisection gives uniqueness of geodesics of length less than $\eps_0$.

{

\begin{wrapfigure}{r}{35mm}
\centering
\vskip-5mm
\includegraphics{mppics/pic-20}
\vskip-0mm
\end{wrapfigure}

As observed above, a pair of points has at most two midpoints.
Suppose midpoints are not always unique, say $p$ and $q$ have two midpoints $m_1$ and $m_2$; set $\eps_0\z=|m_1-m_2|$.

Let $0<|x-y|<\eps_0$.
Choose points $\bar x\z\in [m_1p]$ and $\bar y\in [m_1q]$ at distance $\tfrac12\cdot|x-y|$ from $m_1$.
By the triangle inequality, $|\bar x-\bar y|=|x-y|$.

}

Suppose $x$ and $y$ have two midpoints, say $n_1$ and $n_2$.
By 3-point homogeneity, there is an isometry
$\iota$ sending $x\mapsto \bar x$, $y\mapsto \bar y$, and $n_1\mapsto m_1$.
Then (again by the triangle inequality) the point $m_3=\iota(n_2)$ would be a midpoint of $p$ and $q$.
Note that $m_3\ne m_1$ and $m_3\ne m_2$; the latter follows since
\[|m_3-m_1|\le |x-y|<\eps_0=|m_1-m_2|.\]
Thus the midpoints $m_1$, $m_2$, and $m_3$ are distinct, which is impossible.

\parbf{No branching.}
We will show that \textit{geodesics in $M$ do not branch}, that is, if
two geodesics $[xz]$ and $[xz']$ contain a nontrivial geodesic $[xy]$, then $[xz]$ contains $[xz']$ or $[xz]$ is contained in $[xz']$.

{

\begin{wrapfigure}{r}{35mm}
\centering
\vskip-0mm
\includegraphics{mppics/pic-30}
\vskip-0mm
\end{wrapfigure}

If geodesics branched, we could shorten them and move $x$ closer to the branching point so that $|x-z|\z=|x-z'|\z<\eps_0$ and $z\ne z'$.

By 3-point homogeneity, there is an isometry $\sigma$ fixing $x,y$ and sending $z$ to~$z'$.
But there is also a nonidentity isometry, say $\tau$, fixing $x$ and $z$; it must fix every point on the unique geodesic $[xz]$, including~$y$.
Thus $\sigma$ and $\tau$ are two distinct nonidentity isometries fixing $x$ and $y$ --- a contradiction.

}


\parbf{Lines.}
Any geodesic can be extended to a \emph{local geodesic} $\gamma\colon\mathbb R\to M$; this means that sufficiently short segments of $\gamma$ are geodesics.
For this, use 2-point homogeneity to align a short final segment with the start of a longer geodesic.
By local uniqueness, they agree on the overlap, giving an extension.
Repeating this with the same small step in both directions produces $\gamma$.
The absence of branching ensures that a local geodesic is uniquely determined by any of its nontrivial segments.

Let $d$ be the diameter of $M$; we will show that $d=\infty$.
The following argument is borrowed from Garrett Birkhoff~\cite{birkhoff1944}.

Note that $M$ contains a geodesic of any length less than $d$.
By 2-point homogeneity and the absence of branching, these geodesics can be aligned with the local geodesic $\gamma$ constructed above.
Thus
\[|\gamma(t)-\gamma(s)|=|t-s|\quad\text{if}\quad |t-s|<d.
\eqlbl{eq:|t-s|}\]

Suppose $d<\infty$; then by continuity, \ref{eq:|t-s|} also holds when $|t-s|=d$.
Choose a small $\eps>0$ and set
\[p=\gamma(0),
\quad x=\gamma(d+\eps),
\quad r=|p-x|,
\quad y=\gamma(r),
\quad\text{and}\quad
z=\gamma(\eps).\]
By the triangle inequality, $d-\eps\le r\le d$.

\begin{wrapfigure}{r}{35mm}
\centering
\vskip-8mm
\includegraphics{mppics/pic-40}
\vskip-0mm
\end{wrapfigure}

The points $x$ and $y$ are equidistant from $p$, so there is an isometry $\sigma$ fixing $p$ and swapping $x$ and~$y$.
It reverses the unique geodesic $[xy]$ and hence fixes its midpoint, say $m$.
By our assumptions, there is another nonidentity isometry fixing $p$ and $z$;
it must fix the unique geodesic $[pz]$, and hence every point on $\gamma$.
We now have two distinct nonidentity isometries fixing $p$ and $m$ --- a contradiction.

Thus $d=\infty$, and every segment of the local geodesic $\gamma$ is a geodesic;
that is, $\gamma$ is a geodesic, an isometric embedding of the whole real line into $M$.
Such geodesics will be called \emph{lines}.

Let us check that the geodesic between points $x$ and $y$ is unique.
Suppose there are two.
Extend one of them beyond $y$ to a point $z$.
The other, together with the same segment $[yz]$, is also a geodesic $[xz]$.
The two geodesics branch, which is impossible.

\begin{wrapfigure}{r}{35mm}
\centering
\vskip-5mm
\includegraphics{mppics/pic-50}
\vskip-0mm
\end{wrapfigure}

Thus there is a unique line thru any two distinct points.
There is also a unique nonidentity isometry fixing every point on a given line; we call it the \emph{reflection} in that line.
The assumptions of the theorem also imply that a reflection is an involution; that is, it is its own inverse.

\parbf{Point reflection.}
Choose a point $m$.
Given a point $x$, extend the geodesic $[xm]$ beyond $m$ by its own length.
Denote the opposite endpoint by $x'$ and set $m'=m$.
By uniqueness of the line, the map $\theta\colon x\mapsto x'$ is well defined.
We will show that it is an isometry.
This isometry will be called the \emph{point reflection} about $m$.

\begin{figure}[t!]
\centering
\vskip-0mm
\includegraphics{mppics/pic-60}
\vskip-0mm
\end{figure}

Since $x''=x$, the map $\theta\colon x\mapsto x'$ is bijective, so it suffices to check that
\[|x'-y'|=|x-y|\eqlbl{eq:xy=xy}\]
for any two points $x$ and $y$.
We may assume that $m$, $x$, and $y$ are pairwise distinct; otherwise there is nothing to prove.

Choose $\eps>0$ smaller than $|m-x|$ and $|m-y|$, and points $\bar x\in[mx]$ and $\bar y\in[my]$ at distance $\eps$ from $m$.
By 3-point homogeneity, there is an isometry $\sigma$ fixing $m$ and swapping $\bar x$ and~$\bar y'$.
It sends lines to lines, so it also swaps $\bar x'$ and~$\bar y$.
In particular, $|\bar x-\bar y|=|\bar x'-\bar y'|$.
Again by 3-point homogeneity, there is an isometry fixing $m$ and sending $\bar x\mapsto \bar x'$ and $\bar y\mapsto \bar y'$.
By uniqueness of extension, it sends $x\mapsto x'$ and $y\mapsto y'$;
hence \ref{eq:xy=xy}.

\begin{wrapfigure}{r}{35mm}
\centering
\vskip-8mm
\includegraphics{mppics/pic-100}
\vskip-0mm
\end{wrapfigure}

\parbf{Line reflections.}
Let $\sigma$ be the reflection in a line $\ell$.
We will show that \textit{$\sigma$ fixes only points on $\ell$}.

Choose two distinct points $p$ and $q$ on $\ell$, and let $m$ be their midpoint.
Suppose $\sigma(u)=u$ for some point $u\notin\ell$.
Then $\sigma$ fixes every point on the broken line $[pu]\cup[uq]$.
By uniqueness of lines, the broken line does not pass thru $m$.
As a point $x$ moves along the broken line from $p$ to $q$, the difference $|q-x|-|p-x|$ changes sign.
Thus the broken line contains a point $v$ equidistant from $p$ and $q$.

Let $v'$ be the image of $v$ under point reflection about $m$.
Note that $|p-v|\z=|q-v|=|p-v'|\z=|q-v'|$
and $v'\ne v$, since $v\ne m$.
By 3-point homogeneity, there is an isometry sending $v$ to $v'$ and fixing $p$ and $q$.
By uniqueness, it must be $\sigma$;
that is, $\sigma(v)=v'\ne v$ --- a contradiction.

\parbf{Perpendicular bisector.}
As before, $\sigma$ is the reflection in a line $\ell$.
Let $v\notin\ell$ and $w=\sigma(v)$; we already know that $v\ne w$.

\begin{wrapfigure}{r}{35mm}
\centering
\vskip-0mm
\includegraphics{mppics/pic-110}
\vskip-0mm
\end{wrapfigure}

We will show that
\[x\in \ell\quad\Leftrightarrow\quad|v-x|=|w-x|;\]
the forward implication is obvious.

Let $m$ be the midpoint of $v$ and $w$, and let $\theta$ be the point reflection about $m$.
Note that $\sigma(m)\z=m$, so $m\in\ell$.
Further, $\rho=\theta\circ\sigma$ is the reflection in the line $vw$;
indeed, $\rho$ fixes $v$ and $w$ and reverses~$\ell$.

Suppose $|v-x|=|w-x|$.
If $x\notin \ell$, then $\sigma(x)\ne x$ and $\theta(x)\ne x$.
All three points $x$, $\sigma(x)$, and $\theta(x)$ have the same distances to $v$ and $w$;
thus $\sigma(x)=\theta(x)$ and $x=\rho(x)$, that is, $x$ lies on the line $vw$.
Since $|v-x|=|w-x|$, we get $x=m\in\ell$ --- a contradiction.


\parbf{Half-planes.}
Let
\begin{align*}
V&=\set{x\in M}{|v-x|<|w-x|}
&&\text{and}&
W&=\set{x\in M}{|v-x|>|w-x|}.
\end{align*}
Since $V\ni v$ and $W\ni w$, both sets are nonempty.
A segment joining points in different sets must meet $\ell$,
since the difference of the distances to $v$ and $w$ changes sign along it.

We will show that $V$ and $W$ are \emph{convex}, that is, each contains the geodesic joining any two of its points.

If $x\in V$, then the triangle inequality gives $[vx]\subset V$.
Now let $x,y\in V$.
The broken line $[xv]\cup[vy]$ lies in $V$ and hence does not meet $\ell$.
On this broken line, the distances to $x$ and $\sigma(x)$ cannot be equal; this is possible only on $\ell$.
At its start, $0=|x-x|<|\sigma(x)-x|$, so at its end, $|x-y|<|\sigma(x)-y|$.
If $[xy]$ contained a point $z\in\ell$, then
\[|\sigma(x)-y|\le |\sigma(x)-z|+|z-y|=|x-z|+|z-y|=|x-y|\]
--- a contradiction.
Thus $[xy]$ does not meet $\ell$ and, by continuity, lies entirely in~$V$.
The same argument proves convexity of $W$.

Note that this implies Pasch's axiom: \textit{a line that avoids the vertices of a triangle and meets one of its sides must also meet one of the other two sides.}

\parbf{Final steps.}
Now we need to check that the proved statements imply the axioms of a neutral plane;
once it is done the theorem follows.
This verification is fairly simple, but it depends on the choice of the axiomatic.
We suggest using whichever system of axioms the reader knows best.

For example, one can use Tarski's system~\cite{tarski}, omitting axiom A8 (equivalent to Euclid's fifth postulate).
The notions translate directly into metric terms.
Instead of the axiom schema A13, we can use the stronger full continuity axiom \cite[p. 18]{tarski}; it holds since lines are isometric to $\mathbb R$.
%Axiom A12 follows from the paragraph on the perpendicular bisector.

\parit{We used AI (GPT-6 Astra) to discuss the proof, search the literature, translate, and edit the text.}

{\sloppy
\def\emph{\textit}
\printbibliography[heading=bibintoc]\fussy}
\end{document}

